\chapter{Dari Gen ke Protein}\label{ch:g11:gene-expression}

Pada 1961 seorang ahli biokimia memberi sari bebas \omterm{def:g10:cells-common-unit:cell}{sel} bakteri sebuah
\omterm{def:g11:gene-expression:transcription}{RNA} buatan yang hanya terbuat dari huruf U, yang berulang. Sarinya
membuat \omterm{def:g11:gene-expression:protein}{protein} yang hanya terbuat dari satu \omterm{def:g10:chemistry-of-life:families}{asam amino}, yaitu
fenilalanin, yang berulang. Dengan satu percobaan itu kata pertama \omterm{def:g11:gene-expression:code}{kode
genetiknya} terbaca: UUU berarti fenilalanin. Dalam lima tahun keenam
puluh empat katanya dikenal, dan ternyata semuanya kata yang sama pada
bakteri, tanaman gandum dan manusia. Bab ini mengikuti jalan dari
urutan \omterm{def:g10:universal-dna:nucleotide}{nukleotida} sebuah \omterm{def:g10:universal-dna:gene}{gen} ke urutan \omterm{def:g10:chemistry-of-life:families}{asam amino} sebuah \omterm{def:g11:gene-expression:protein}{protein} ---
jalan yang dilalui tiap sifat pada \cref{ch:g11:mutations}.

\section{Gen membuat protein}

\begin{proposition}[Satu gen, satu protein]\label{prop:g11:gene-expression:onegene}
Sebuah \omterm{def:g10:universal-dna:gene}{gen} terungkap ketika \omterm{def:g10:cells-common-unit:cell}{selnya} memakai urutannya untuk membuat
\omterm{def:g11:gene-expression:protein}{protein} yang disandikannya. Sebagian besar sifat bergantung pada
\omterm{def:g11:gene-expression:protein}{protein} --- \omterm{def:g10:cell-metabolism:metabolism}{enzim}, serat struktur, pengangkut, reseptor --- dan sebuah
\omterm{def:g11:mutations:mutation}{mutasi} mengubah sifat dengan mengubah \omterm{def:g11:gene-expression:protein}{protein} yang dihasilkan \omterm{def:g10:universal-dna:gene}{gennya}.
\end{proposition}
\begin{proof}[Bukti percobaan]
Beadle dan Tatum (1941) menyinari spora kapang roti dan mengumpulkan
mutan yang tidak lagi dapat tumbuh pada medium minimal kecuali jika
satu zat tertentu, seperti \omterm{def:g10:chemistry-of-life:families}{asam amino} arginin, ditambahkan. Tiap
mutannya kehilangan satu \omterm{def:g10:cell-metabolism:metabolism}{enzim} dalam rantai reaksi yang membuat zat
itu; tiap cacatnya diwariskan sebagai satu \omterm{def:g10:universal-dna:gene}{gen}; dan mutan yang
terhalang pada langkah yang berbeda dalam rantai yang sama membawa
\omterm{def:g11:mutations:mutation}{mutasi} pada \omterm{def:g10:universal-dna:gene}{gen} yang berbeda. Satu \omterm{def:g10:universal-dna:gene}{gen}, satu \omterm{def:g10:cell-metabolism:metabolism}{enzim} --- dan, sebagaimana
digeneralkan penelitian kemudian, satu \omterm{def:g10:universal-dna:gene}{gen}, satu rantai polipeptida.
\end{proof}

\begin{definition}[Protein, urutan asam amino]\label{def:g11:gene-expression:protein}
Sebuah \emph{protein}\index{protein} adalah rantai \omterm{def:g10:chemistry-of-life:families}{asam amino},
panjangnya dari beberapa lusin sampai beberapa ribu, yang diambil dari
perangkat dua puluh jenis. \emph{Urutan} --- yaitu susunan \omterm{def:g10:chemistry-of-life:families}{asam
aminonya} --- adalah apa yang ditetapkan \omterm{def:g10:universal-dna:gene}{gennya}; begitu terbuat,
rantainya melipat menjadi bentuk tiga dimensi yang tertentu yang
ditetapkan urutan itu, dan bentuknya menetapkan fungsinya. Ubahlah satu
\omterm{def:g10:chemistry-of-life:families}{asam amino} dan lipatannya, serta fungsinya, dapat berubah.
\end{definition}

\section{Transkripsi: gennya disalin menjadi RNA}

\begin{definition}[RNA duta dan transkripsi]\label{def:g11:gene-expression:transcription}
\emph{RNA}\index{RNA} (asam ribonukleat) adalah rantai \omterm{def:g10:universal-dna:nucleotide}{nukleotida}
berutai tunggal seperti \omterm{def:g10:universal-dna:nucleotide}{nukleotida} \omterm{def:g10:universal-dna:information}{DNA}, kecuali bahwa gulanya adalah
ribosa dan \omterm{def:g10:universal-dna:nucleotide}{basa} urasil (U) menggantikan timin (T).
\emph{Transkripsi}\index{transkripsi} adalah penyalinan urutan sebuah
\omterm{def:g10:universal-dna:gene}{gen} menjadi molekul \emph{RNA duta}\index{RNA duta} (mRNA): \omterm{def:g10:cell-metabolism:metabolism}{enzim}
\emph{RNA polimerase} membuka \omterm{prop:g10:universal-dna:helix}{heliks gandanya} di sepanjang \omterm{def:g10:universal-dna:gene}{gennya},
membaca satu untai (yaitu \emph{untai cetakan}) dan merakit RNA
komplementernya, U berhadapan dengan A, A berhadapan dengan T, G
berhadapan dengan C dan C berhadapan dengan G. RNA duta itu punya
urutan untai \omterm{def:g10:universal-dna:gene}{gen} yang lain, dengan U untuk T; ia meninggalkan \omterm{def:g10:cells-common-unit:organelle}{inti
selnya} menuju \omterm{def:g10:cells-common-unit:cell}{sitoplasma}.
\end{definition}

\begin{omfigure}
\begin{tikzpicture}[font=\small, >=Stealth, line cap=round]
  % DNA strands with bubble
  \draw[very thick, omDef] (0,0.5) -- (2.5,0.5) .. controls (3.2,0.5) and (3.4,1.3) .. (4.3,1.3) -- (6.7,1.3) .. controls (7.6,1.3) and (7.8,0.5) .. (8.5,0.5) -- (11,0.5);
  \draw[very thick, omDef] (0,-0.5) -- (2.5,-0.5) .. controls (3.2,-0.5) and (3.4,-1.3) .. (4.3,-1.3) -- (6.7,-1.3) .. controls (7.6,-1.3) and (7.8,-0.5) .. (8.5,-0.5) -- (11,-0.5);
  \foreach \x in {0.3,0.7,...,2.3} \draw[gray] (\x,0.5) -- (\x,-0.5);
  \foreach \x in {8.7,9.1,...,10.9} \draw[gray] (\x,0.5) -- (\x,-0.5);
  % template strand letters (bottom strand in bubble)
  \foreach \x/\l in {4.5/T, 4.9/A, 5.3/C, 5.7/G, 6.1/G, 6.5/A} \node[font=\scriptsize, omDef, anchor=north] at (\x,-1.32) {\l};
  % mRNA
  \draw[very thick, omThm] (4.4,-0.9) -- (6.6,-0.9) .. controls (7.4,-0.9) and (7.2,0.1) .. (8.0,0.6) .. controls (8.6,1.0) and (9.4,1.4) .. (10.6,1.7);
  \foreach \x/\l in {4.5/A, 4.9/U, 5.3/G, 5.7/C, 6.1/C, 6.5/U} \node[font=\scriptsize, omThm, anchor=south] at (\x,-0.85) {\l};
  \foreach \x in {4.5,4.9,...,6.5} \draw[gray, densely dotted] (\x,-0.95) -- (\x,-1.28);
  % polymerase
  \draw[thick, fill=omProp!30, rounded corners=8pt] (6.95,-1.6) rectangle (8.65,0.2);
  \node[align=center] at (7.8,-0.7) {RNA\\ polimerase};
  \draw[->, thick] (8.7,-1.9) -- (9.8,-1.9) node[right] {bergerak sepanjang gennya};
  \node[anchor=east, omDef, align=right] at (-0.15,0) {DNA\\ gennya};
  \node[anchor=west, omThm, align=left] at (10.7,1.7) {mRNA, yang sedang\\ dilepaskan};
  \node[anchor=north, omDef] at (5.5,-1.75) {untai cetakan};
\end{tikzpicture}

{\small \omterm{def:g11:gene-expression:transcription}{Transkripsi}. \omterm{def:g11:gene-expression:transcription}{RNA} polimerase membuka heliksnya, membaca untai
cetakannya dan membangun sebuah \omterm{def:g11:gene-expression:transcription}{RNA duta} yang komplementer dengannya (U
sebagai ganti T). Di belakang \omterm{def:g10:cell-metabolism:metabolism}{enzimnya} heliksnya menutup lagi; \omterm{def:g11:gene-expression:transcription}{RNA-nya}
mengelupas.}
\end{omfigure}

\begin{example}[Sebuah transkrip]\label{ex:g11:gene-expression:transcript}
Untai cetakan \texttt{3'-TAC GGA CTT ATC-5'} memberi
mRNA \texttt{5'-AUG CCU GAA UAG-3'}. Untai \omterm{def:g10:universal-dna:gene}{gen} yang lain berbunyi
\texttt{5'-ATG CCT GAA TAG-3'}: mRNA-nya adalah urutan untai itu
dengan U untuk T, dan itulah sebabnya urutan \omterm{def:g10:universal-dna:gene}{gen} menurut kesepakatan
dituliskan sebagai untai itu, yaitu \emph{untai penyandi}.
\end{example}

\section{Kode genetiknya}

\begin{definition}[Kodon dan kode genetik]\label{def:g11:gene-expression:code}
\omterm{def:g11:gene-expression:transcription}{RNA dutanya} dibaca dalam kelompok tiga \omterm{def:g10:universal-dna:nucleotide}{nukleotida} yang berurutan, yaitu
\emph{kodon}\index{kodon}, dari sebuah titik awal yang tertentu. Yang
disebut \emph{kode genetik}\index{kode genetik} adalah kesepadanan
antara 64 kodon yang mungkin dan 20 \omterm{def:g10:chemistry-of-life:families}{asam aminonya}: 61 kodon menetapkan
satu \omterm{def:g10:chemistry-of-life:families}{asam amino}, tiga kodon (UAA, UAG, UGA) merupakan isyarat
\emph{henti} yang mengakhiri rantainya, dan AUG, yang menetapkan
metionin, juga bertugas sebagai kodon \emph{awal}. Kodenya bersifat
\emph{degenerasi} --- sebagian besar \omterm{def:g10:chemistry-of-life:families}{asam aminonya} punya lebih dari
satu kodon --- dan \emph{semesta}: tabel yang sama melayani setiap organisme yang
dikenal, dengan sedikit ragam kecil yang langka.
\end{definition}

\begin{omfigure}
\begin{tabular}{c|cccc|c}
\multicolumn{1}{c}{} & \multicolumn{4}{c}{\small huruf kedua} & \multicolumn{1}{c}{} \\
{\small pertama} & U & C & A & G & {\small ketiga} \\ \hline
U & UUU Phe & UCU Ser & UAU Tyr & UGU Cys & U \\
 & UUC Phe & UCC Ser & UAC Tyr & UGC Cys & C \\
 & UUA Leu & UCA Ser & UAA \textbf{henti} & UGA \textbf{henti} & A \\
 & UUG Leu & UCG Ser & UAG \textbf{henti} & UGG Trp & G \\ \hline
C & CUU Leu & CCU Pro & CAU His & CGU Arg & U \\
 & CUC Leu & CCC Pro & CAC His & CGC Arg & C \\
 & CUA Leu & CCA Pro & CAA Gln & CGA Arg & A \\
 & CUG Leu & CCG Pro & CAG Gln & CGG Arg & G \\ \hline
A & AUU Ile & ACU Thr & AAU Asn & AGU Ser & U \\
 & AUC Ile & ACC Thr & AAC Asn & AGC Ser & C \\
 & AUA Ile & ACA Thr & AAA Lys & AGA Arg & A \\
 & AUG \textbf{Met} & ACG Thr & AAG Lys & AGG Arg & G \\ \hline
G & GUU Val & GCU Ala & GAU Asp & GGU Gly & U \\
 & GUC Val & GCC Ala & GAC Asp & GGC Gly & C \\
 & GUA Val & GCA Ala & GAA Glu & GGA Gly & A \\
 & GUG Val & GCG Ala & GAG Glu & GGG Gly & G \\
\end{tabular}

{\small \omterm{def:g11:gene-expression:code}{Kode genetik}, yang dibaca pada mRNA-nya. Tiap \omterm{def:g11:gene-expression:code}{kodon} ditemukan
lewat huruf pertamanya (baris), huruf keduanya (lajur) dan huruf
ketiganya (baris di dalam bloknya). AUG sekaligus metionin dan isyarat
awal; tiga \omterm{def:g11:gene-expression:code}{kodon} merupakan henti.}
\end{omfigure}

\begin{proposition}[Bagaimana kodenya terbaca]\label{prop:g11:gene-expression:readcode}
Makna tiap \omterm{def:g11:gene-expression:code}{kodon} ditegakkan lewat percobaan.
\end{proposition}
\begin{proof}[Bukti percobaan]
Nirenberg dan Matthaei (1961) menambahkan \omterm{def:g11:gene-expression:transcription}{RNA} sintetis dari satu
\omterm{def:g10:universal-dna:nucleotide}{nukleotida} yang berulang ke dalam sari bebas \omterm{def:g10:cells-common-unit:cell}{sel} yang memuat \omterm{def:g10:cells-common-unit:organelle}{ribosom},
\omterm{def:g10:chemistry-of-life:families}{asam amino} dan \omterm{def:g10:cell-metabolism:metabolism}{enzim} penyusun \omterm{def:g11:gene-expression:protein}{protein}: poli-U menghasilkan rantai
fenilalanin, poli-C rantai prolin, poli-A rantai lisin. \omterm{def:g11:gene-expression:transcription}{RNA} sintetis
dari pasangan dan tripel yang berulang, dan kemudian pengikatan tripel
tunggal pada \omterm{def:g10:cells-common-unit:organelle}{ribosom}, menetapkan \omterm{def:g11:gene-expression:code}{kodon} selebihnya pada 1966. Panjang
tiga hurufnya telah disimpulkan dari hitungannya (dua huruf hanya
memberi 16 gabungan, terlalu sedikit untuk 20 \omterm{def:g10:chemistry-of-life:families}{asam amino}; tiga huruf
memberi 64) dan dibenarkan oleh \omterm{def:g11:mutations:mutation}{mutasi}: menyisipkan satu atau dua
\omterm{def:g10:universal-dna:nucleotide}{nukleotida} ke dalam sebuah \omterm{def:g10:universal-dna:gene}{gen} memusnahkan \omterm{def:g11:gene-expression:protein}{proteinnya}, menyisipkan tiga
memulihkan \omterm{def:g11:gene-expression:protein}{protein} yang nyaris normal.
\end{proof}

\begin{method}[Menerjemahkan sebuah urutan]\label{met:g11:gene-expression:translate}
\begin{enumerate}
  \item Tuliskan untai penyandinya (atau transkripsikan untai
        cetakannya): gantilah T dengan U untuk memperoleh mRNA-nya.
  \item Carilah AUG yang pertama: itulah awalnya, dan rangka bacanya
        tertetapkan darinya.
  \item Potonglah mRNA-nya menjadi tripel yang berurutan mulai dari
        AUG-nya lalu bacalah masing-masing pada tabelnya.
  \item Berhentilah pada \omterm{def:g11:gene-expression:code}{kodon} henti yang pertama; \omterm{def:g10:chemistry-of-life:families}{asam amino} yang
        terdaftar, berurutan, adalah urutan \omterm{def:g11:gene-expression:protein}{proteinnya}.
  \item Untuk sebuah mutan, ulangi lalu bandingkan: \omterm{def:g11:gene-expression:protein}{protein} yang sama
        (diam), satu \omterm{def:g10:chemistry-of-life:families}{asam amino} yang berubah (salah makna), henti yang
        terlalu dini (tanpa makna), segalanya berubah setelah \omterm{def:g11:mutations:mutation}{mutasinya}
        (pergeseran rangka).
\end{enumerate}
\end{method}

\begin{example}[Empat mutasi, empat hasil]\label{ex:g11:gene-expression:outcomes}
Untai penyandi \texttt{ATG CCT GAA TTC TAG}: mRNA-nya
\texttt{AUG CCU GAA UUC UAG}, \omterm{def:g11:gene-expression:protein}{proteinnya} Met--Pro--Glu--Phe.
\begin{itemize}
  \item \texttt{ATG CC\textbf{C} GAA TTC TAG}: CCC tetap Pro ---
        \emph{diam}.
  \item \texttt{ATG CCT G\textbf{C}A TTC TAG}: GCA adalah Ala ---
        Met--Pro--Ala--Phe, \emph{salah makna}.
  \item \texttt{ATG CCT \textbf{T}AA TTC TAG}: UAA adalah henti ---
        Met--Pro, \emph{tanpa makna}, sebuah \omterm{def:g11:gene-expression:protein}{protein} yang terpenggal.
  \item \texttt{ATG CC\textbf{A}T GAA TTC TAG}: A yang tersisip
        menggeser rangkanya: \texttt{AUG CCA UGA}\dots{} Met--Pro--henti
        --- \emph{pergeseran rangka}.
\end{itemize}
\end{example}

\section{Translasi: pesannya dibaca}

\begin{proposition}[Translasi]\label{prop:g11:gene-expression:translation}
\emph{Translasi}\index{translasi} berlangsung pada
\emph{\omterm{def:g10:cells-common-unit:organelle}{ribosom}}, di dalam \omterm{def:g10:cells-common-unit:cell}{sitoplasma}. Sebuah \omterm{def:g10:cells-common-unit:organelle}{ribosom} mengikat mRNA-nya
pada \omterm{def:g11:gene-expression:code}{kodon} awalnya dan bergerak sepanjangnya \omterm{def:g11:gene-expression:code}{kodon} demi \omterm{def:g11:gene-expression:code}{kodon}. Tiap
\omterm{def:g11:gene-expression:code}{kodon} dipasangkan oleh sebuah \emph{\omterm{def:g11:gene-expression:transcription}{RNA} transfer} (tRNA), yaitu \omterm{def:g11:gene-expression:transcription}{RNA}
kecil yang membawa, pada satu ujungnya, \emph{antikodon} bernukleotida
tiga yang komplementer dengan \omterm{def:g11:gene-expression:code}{kodonnya} dan, pada ujung yang lain, \omterm{def:g10:chemistry-of-life:families}{asam
amino} yang sepadan. \omterm{def:g10:cells-common-unit:organelle}{Ribosomnya} menyambungkan tiap \omterm{def:g10:chemistry-of-life:families}{asam amino} yang
dibawa masuk ke rantai yang sedang tumbuh dan melepaskan tRNA yang
kosong; pada \omterm{def:g11:gene-expression:code}{kodon} henti ia melepaskan rantai yang selesai, yang lalu
melipat. Beberapa \omterm{def:g10:cells-common-unit:organelle}{ribosom} membaca satu mRNA berturut-turut, dan satu
mRNA menghasilkan ratusan salinan \omterm{def:g11:gene-expression:protein}{proteinnya} sebelum ia diuraikan.
\end{proposition}
\admitted

\begin{omfigure}
\begin{tikzpicture}[font=\small, >=Stealth, line cap=round]
  % mRNA
  \draw[very thick, omThm] (0,0) -- (11,0);
  \foreach \x/\l in {0.4/A,0.8/U,1.2/G, 1.7/C,2.1/C,2.5/U, 3.0/G,3.4/A,3.8/A, 4.3/U,4.7/U,5.1/C, 5.6/G,6.0/G,6.4/U, 6.9/U,7.3/A,7.7/G}
    \node[font=\scriptsize, omThm, anchor=north] at (\x,-0.05) {\l};
  \node[anchor=east, omThm] at (-0.15,0) {mRNA};
  % ribosome (two lobes) over codons 3 and 4
  \draw[thick, fill=omProp!25, rounded corners=10pt] (2.7,0.15) rectangle (5.5,1.4);
  \draw[thick, fill=omProp!35, rounded corners=12pt] (2.5,1.4) .. controls (2.5,3.0) and (5.7,3.0) .. (5.7,1.4) -- cycle;
  \node at (4.1,2.1) {ribosom};
  % tRNAs
  \draw[very thick, omMeth!80!black] (3.4,0.55) -- (3.4,1.35);
  \node[font=\scriptsize, omMeth!80!black, anchor=south] at (3.4,0.15) {CUU};
  \draw[very thick, omMeth!80!black] (4.7,0.55) -- (4.7,1.35);
  \node[font=\scriptsize, omMeth!80!black, anchor=south] at (4.7,0.15) {AAG};
  % amino acids chain
  \foreach \i/\x in {0/3.4, 1/4.7} \fill[omDef!70] (\x,1.5) circle (4pt);
  \foreach \x in {2.6, 1.9, 1.2} \fill[omDef!70] (\x,1.9) circle (4pt);
  \draw[thick, omDef!70] (1.2,1.9) -- (2.6,1.9) -- (3.4,1.5) -- (4.7,1.5);
  \node[anchor=east, omDef!80!black, align=right] at (0.9,1.9) {protein yang\\ sedang tumbuh};
  % incoming tRNA
  \draw[very thick, omMeth!80!black] (7.0,1.3) -- (7.0,2.1);
  \fill[omDef!70] (7.0,2.25) circle (4pt);
  \node[font=\scriptsize, omMeth!80!black, anchor=west] at (7.15,1.35) {CCA};
  \node[anchor=west, align=left, font=\scriptsize] at (7.3,1.7) {tRNA berikutnya, dengan\\ asam aminonya, tiba};
  \draw[->, thick] (7.0,1.1) -- (6.1,0.75);
  \draw[->, thick] (5.7,-0.9) -- (7.2,-0.9) node[right] {ribosomnya bergerak};
  \node[anchor=north, font=\scriptsize] at (4.05,-0.45) {kodon};
  \draw (2.85,-0.35) -- (2.85,-0.45) -- (3.95,-0.45) -- (3.95,-0.35);
\end{tikzpicture}

{\small \omterm{prop:g11:gene-expression:translation}{Translasi}. \omterm{def:g10:cells-common-unit:organelle}{Ribosomnya} menahan dua \omterm{def:g11:gene-expression:code}{kodon}; sebuah tRNA yang
antikodonnya cocok dengan tiap \omterm{def:g11:gene-expression:code}{kodonnya} membawa \omterm{def:g10:chemistry-of-life:families}{asam aminonya},
rantainya dipindahkan ke pendatang barunya, dan \omterm{def:g10:cells-common-unit:organelle}{ribosomnya} melangkah
satu \omterm{def:g11:gene-expression:code}{kodon}. Pada \omterm{def:g11:gene-expression:code}{kodon} henti rantainya dilepaskan.}
\end{omfigure}

\begin{example}[Kecepatan dan hasilnya]\label{ex:g11:gene-expression:speed}
Sebuah \omterm{def:g10:cells-common-unit:organelle}{ribosom} menambahkan 5 sampai 20 \omterm{def:g10:chemistry-of-life:families}{asam amino} tiap sekon: \omterm{def:g11:gene-expression:protein}{protein}
berisi 300 \omterm{def:g10:chemistry-of-life:families}{asam amino} memakan kurang dari satu menit. \omterm{def:g10:cells-common-unit:organelle}{Ribosom} saling
menyusul pada sebuah mRNA berjarak sekitar 80 \omterm{def:g10:universal-dna:nucleotide}{nukleotida}, sehingga
pesan sepanjang 900 \omterm{def:g10:universal-dna:nucleotide}{nukleotida} membawa selusin \omterm{def:g10:cells-common-unit:organelle}{ribosom} sekaligus;
sepanjang hidupnya yang beberapa jam, satu mRNA menghasilkan seribu
molekul \omterm{def:g11:gene-expression:protein}{protein}. \omterm{def:g10:cells-common-unit:cell}{Sel} \emph{E.~coli} membuat sekitar $\num{20000}$
molekul \omterm{def:g11:gene-expression:protein}{protein} tiap sekon; bakal \omterm{def:g10:cells-common-unit:cell}{sel} darah merah, sekitar
$\num{5e8}$ molekul hemoglobin seumur hidupnya.
\end{example}

\section{Satu gen, beberapa pesan}

\begin{proposition}[Pematangan pesan pada eukariot]\label{prop:g11:gene-expression:splicing}
Pada \omterm{def:g10:cells-common-unit:prokaryote}{sel eukariotik} urutan sebuah \omterm{def:g10:universal-dna:gene}{gen} diselingi
\emph{intron}\index{intron}, yaitu ruas yang ditranskripsikan tetapi
lalu dipotong keluar dari \omterm{def:g11:gene-expression:transcription}{RNA-nya} sebelum \omterm{def:g11:gene-expression:transcription}{RNA} itu meninggalkan \omterm{def:g10:cells-common-unit:organelle}{inti
selnya}; ruas yang disimpan dan disambungkan ujung ke ujung, yaitu
\emph{ekson}, membentuk mRNA yang matang. Banyak \omterm{def:g10:universal-dna:gene}{gen} dapat dipotong
dengan lebih dari satu cara, dengan menyimpan perangkat ekson yang
berbeda (\emph{penyambungan alternatif}), sehingga satu \omterm{def:g10:universal-dna:gene}{gen} menghasilkan
beberapa mRNA yang berbeda dan beberapa \omterm{def:g11:gene-expression:protein}{protein} yang berkerabat:
ke-$\num{20000}$ \omterm{def:g10:universal-dna:gene}{gen} manusia menyandikan jauh lebih dari
$\num{100000}$ \omterm{def:g11:gene-expression:protein}{protein}.
\end{proposition}
\admitted

\begin{omfigure}
\begin{tikzpicture}[font=\small, >=Stealth]
  \tikzset{ex/.style={draw, fill=omDef!40, minimum height=0.45cm, minimum width=1.0cm, inner sep=1pt},
           intr/.style={draw, fill=gray!20, minimum height=0.45cm, minimum width=0.7cm, inner sep=1pt}}
  \node[anchor=east] at (-0.2,0) {gen (DNA)};
  \node[ex] at (0.5,0) {1}; \node[intr] at (1.35,0) {}; \node[ex] at (2.2,0) {2}; \node[intr] at (3.05,0) {}; \node[ex] at (3.9,0) {3}; \node[intr] at (4.75,0) {}; \node[ex] at (5.6,0) {4};
  \node[anchor=west, font=\scriptsize] at (6.2,0) {ekson (biru), intron (kelabu)};
  \draw[->, thick] (3.0,-0.35) -- (3.0,-0.95) node[midway, right] {transkripsi};
  \node[anchor=east] at (-0.2,-1.3) {pra-mRNA};
  \node[ex, fill=omThm!30] at (0.5,-1.3) {1}; \node[intr] at (1.35,-1.3) {}; \node[ex, fill=omThm!30] at (2.2,-1.3) {2}; \node[intr] at (3.05,-1.3) {}; \node[ex, fill=omThm!30] at (3.9,-1.3) {3}; \node[intr] at (4.75,-1.3) {}; \node[ex, fill=omThm!30] at (5.6,-1.3) {4};
  \draw[->, thick] (2.0,-1.65) -- (1.0,-2.35) node[midway, left] {penyambungan};
  \draw[->, thick] (4.0,-1.65) -- (6.0,-2.35) node[midway, right] {penyambungan alternatif};
  \node[anchor=east] at (-0.6,-2.7) {mRNA A};
  \node[ex, fill=omThm!30] at (0.0,-2.7) {1}; \node[ex, fill=omThm!30] at (1.0,-2.7) {2}; \node[ex, fill=omThm!30] at (2.0,-2.7) {3}; \node[ex, fill=omThm!30] at (3.0,-2.7) {4};
  \node[anchor=west] at (3.6,-2.7) {protein A};
  \node[anchor=east] at (5.4,-2.7) {};
  \node[ex, fill=omThm!30] at (6.0,-2.7) {1}; \node[ex, fill=omThm!30] at (7.0,-2.7) {2}; \node[ex, fill=omThm!30] at (8.0,-2.7) {4};
  \node[anchor=west] at (8.6,-2.7) {protein B};
\end{tikzpicture}

{\small Sebuah \omterm{def:g10:universal-dna:gene}{gen} eukariotik ditranskripsikan utuh, lalu \omterm{prop:g11:gene-expression:splicing}{intronnya}
disingkirkan. Menyimpan seluruh eksonnya, atau meninggalkan satu di
antaranya, memberi dua mRNA dan dua \omterm{def:g11:gene-expression:protein}{protein} dari \omterm{def:g10:universal-dna:gene}{gen} yang sama.}
\end{omfigure}

\begin{remark}[Aliran keterangannya]\label{rem:g11:gene-expression:flow}
\omterm{def:g10:universal-dna:information}{DNA} ditranskripsikan menjadi \omterm{def:g11:gene-expression:transcription}{RNA}, \omterm{def:g11:gene-expression:transcription}{RNA} diterjemahkan menjadi \omterm{def:g11:gene-expression:protein}{protein},
dan \omterm{def:g11:gene-expression:protein}{protein} membuat sifatnya: keterangannya mengalir satu arah. \omterm{def:g11:gene-expression:protein}{Protein}
yang berubah tidak pernah menulis ulang \omterm{def:g10:universal-dna:gene}{gennya}, dan itulah sebabnya
latihan atau kulit terbakar matahari seumur hidup tidak diwariskan, dan
sebabnya hanya \omterm{def:g11:mutations:mutation}{mutasi} \omterm{def:g10:universal-dna:information}{DNA} (\cref{ch:g11:mutations}) yang mengubah apa
yang diterima keturunan berikutnya. Ketiga langkah yang sama berjalan
dalam setiap \omterm{def:g10:cells-common-unit:cell}{sel} setiap organisme, dalam kode yang sama: yaitu
kesemestaan \cref{ch:g10:universal-dna}, kini sampai ke kata terakhirnya.
\end{remark}

\section{Latihan}

\begin{exercise}[$\star$]\label{exo:g11:gene-expression:1}
Berikan tiga beda antara \omterm{def:g10:universal-dna:information}{DNA} dan \omterm{def:g11:gene-expression:transcription}{RNA}.
\end{exercise}

\begin{exercise}[$\star$]\label{exo:g11:gene-expression:2}
Transkripsikan untai cetakan \texttt{3'-TAC CGA AAA ATT-5'} menjadi
mRNA.
\end{exercise}

\begin{exercise}[$\star$]\label{exo:g11:gene-expression:3}
Terjemahkan mRNA \texttt{AUG GGC AAA UGU UAA} dengan tabel kodenya.
\end{exercise}

\begin{exercise}[$\star$]\label{exo:g11:gene-expression:4}
Apa peran \omterm{def:g10:cells-common-unit:organelle}{ribosom} dan peran \omterm{def:g11:gene-expression:transcription}{RNA} transfer dalam translasi?
\end{exercise}

\begin{exercise}[$\star$]\label{exo:g11:gene-expression:5}
Mengapa kodenya harus memakai sedikitnya tiga \omterm{def:g10:universal-dna:nucleotide}{nukleotida} tiap \omterm{def:g10:chemistry-of-life:families}{asam amino}?
\end{exercise}

\begin{exercise}[$\star\star$]\label{exo:g11:gene-expression:6}
Sebuah \omterm{def:g11:gene-expression:protein}{protein} punya 412 \omterm{def:g10:chemistry-of-life:families}{asam amino}. Berapa panjang paling sedikit
bagian penyandi mRNA-nya, termasuk \omterm{def:g11:gene-expression:code}{kodon} hentinya, dan berapa panjang
\omterm{def:g10:universal-dna:gene}{gen} yang sepadan dalam pasangan \omterm{def:g10:universal-dna:nucleotide}{basa}?
\end{exercise}

\begin{exercise}[$\star\star$]\label{exo:g11:gene-expression:7}
Untai penyandi \texttt{ATG AAA GGC TGG TAA} bermutasi menjadi
\texttt{ATG AAA GGC TGA TAA}. Berikan kedua \omterm{def:g11:gene-expression:protein}{proteinnya} dan golongkan
\omterm{def:g11:mutations:mutation}{mutasinya}.
\end{exercise}

\begin{exercise}[$\star\star$]\label{exo:g11:gene-expression:8}
Urutan awal yang sama, bermutasi menjadi \texttt{ATG AAG GGC TGG TAA}
dan menjadi \texttt{ATG AAA GGG CTG GTA A}. Golongkan masing-masing dan
berikan \omterm{def:g11:gene-expression:protein}{proteinnya}.
\end{exercise}

\begin{exercise}[$\star\star$]\label{exo:g11:gene-expression:9}
Jelaskan mengapa substitusi pada kedudukan ketiga sebuah \omterm{def:g11:gene-expression:code}{kodon} sering
bersifat diam, dengan memakai tabelnya.
\end{exercise}

\begin{exercise}[$\star\star$]\label{exo:g11:gene-expression:10}
\omterm{def:g11:gene-expression:transcription}{RNA} poli-UC (\texttt{UCUCUCUC}\dots) memberi \omterm{def:g11:gene-expression:protein}{protein} yang berselang-seling
serin dan leusin. Tunjukkan bahwa ini taat asas dengan kode tiga huruf
dan pakailah untuk menetapkan dua \omterm{def:g11:gene-expression:code}{kodon}.
\end{exercise}

\begin{exercise}[$\star\star$]\label{exo:g11:gene-expression:11}
Sebuah \omterm{def:g10:universal-dna:gene}{gen} manusia sepanjang \num{30000} pasangan \omterm{def:g10:universal-dna:nucleotide}{basa} menghasilkan
\omterm{def:g11:gene-expression:protein}{protein} berisi 500 \omterm{def:g10:chemistry-of-life:families}{asam amino}. Jelaskan ketimpangan itu.
\end{exercise}

\begin{exercise}[$\star\star\star$]\label{exo:g11:gene-expression:12}
Sebuah obat menghalangi \omterm{def:g10:cells-common-unit:organelle}{ribosom} bakteri tetapi tidak menghalangi
\omterm{def:g10:cells-common-unit:organelle}{ribosom} manusia. Jelaskan mengapa ia dapat menjadi antibiotik, dan apa
makna keberadaannya tentang \omterm{def:g10:cells-common-unit:organelle}{ribosom} pada kedua kelompok itu.
\end{exercise}

\begin{exercise}[$\star\star\star$]\label{exo:g11:gene-expression:13}
Jelaskan bagaimana \omterm{def:g10:universal-dna:gene}{gen} ubur-ubur pada \cref{ch:g10:universal-dna} dapat
dibaca seekor mencit, dalam istilah bab ini, dan apa yang akan terjadi
jika mencitnya memakai kode yang berbeda.
\end{exercise}

\begin{exercise}[$\star\star\star$]\label{exo:g11:gene-expression:14}
Sebuah \omterm{def:g11:mutations:mutation}{mutasi} mengubah tRNA yang antikodonnya membaca UAG sehingga ia
membawa sebuah \omterm{def:g10:chemistry-of-life:families}{asam amino} alih-alih berhenti. Ramalkan pengaruhnya pada
\omterm{def:g11:gene-expression:protein}{protein} \omterm{def:g10:cells-common-unit:cell}{selnya} secara umum, dan pada \omterm{def:g10:universal-dna:gene}{gen} mutan yang membawa UAG yang
terlalu dini.
\end{exercise}

\begin{exercise}[$\star\star\star$]\label{exo:g11:gene-expression:15}
Bahaslah mengapa pergeseran rangka di dekat awal sebuah \omterm{def:g10:universal-dna:gene}{gen} hampir
selalu lebih buruk daripada pergeseran di dekat ujungnya, dan mengapa
\omterm{def:g11:mutations:mutation}{mutasi} salah makna dapat berkisar dari tak berbahaya sampai mematikan.
\end{exercise}

\section{Soal: Sebuah Gen yang Dibaca Empat Cara}

\begin{problem}[{Soal akhir pekan --- satu gen pendek yang
ditranskripsikan, diterjemahkan, dimutasikan tiga kali dan diukur
waktunya: dari DNA-nya ke asam aminonya, dan enam puluh sekon yang
dituntut sebuah protein}]
\label{pb:g11:gene-expression:1}
Untai penyandi sebuah \omterm{def:g10:universal-dna:gene}{gen} pendek berbunyi:
\begin{center}
\texttt{ATG GCT TGG AAA CCC
GAA TAC GGT CAT TTC TGA}
\end{center}

\medskip
\noindent\textbf{Bagian I --- Membaca \omterm{def:g10:universal-dna:gene}{gennya}.}
\begin{enumerate}
  \item Tuliskan untai cetakannya.
  \item Tuliskan mRNA yang ditranskripsikan dari \omterm{def:g10:universal-dna:gene}{gennya}.
  \item Terjemahkan mRNA itu, \omterm{def:g11:gene-expression:code}{kodon} demi \omterm{def:g11:gene-expression:code}{kodon}, dan berikan urutan \omterm{def:g10:chemistry-of-life:families}{asam
        amino} \omterm{def:g11:gene-expression:protein}{proteinnya}.
  \item Berapa banyak \omterm{def:g10:chemistry-of-life:families}{asam amino} yang dimuat \omterm{def:g11:gene-expression:protein}{proteinnya}, dan berapa
        banyak \omterm{def:g10:universal-dna:nucleotide}{nukleotida} mRNA-nya yang dipakai untuk menetapkannya,
        termasuk hentinya?
  \item \omterm{def:g11:gene-expression:code}{Kodon} mana yang memulai pembacaannya, dan apa yang menetapkan
        rangka baca bagi semua \omterm{def:g11:gene-expression:code}{kodon} lainnya?
\end{enumerate}

\medskip
\noindent\textbf{Bagian II --- Tiga mutan.}
\begin{enumerate}[resume]
  \item Mutan 1: \texttt{ATG GCT TGG AAA CCG
        GAA TAC GGT CAT TTC TGA}. Berikan \omterm{def:g11:gene-expression:protein}{proteinnya} dan golongkan \omterm{def:g11:mutations:mutation}{mutasinya}.
  \item Mutan 2: \texttt{ATG GCT TGG AAA CCC
        GAA TAC GGT CAT TTC TGA} dengan \omterm{def:g11:gene-expression:code}{kodon} kesembilannya
        \texttt{CAT} digantikan \texttt{CGT}.
        Berikan \omterm{def:g11:gene-expression:protein}{proteinnya} dan golongkan.
  \item Mutan 3: \texttt{ATG GCT TGA AAA CCC
        GAA TAC GGT CAT TTC TGA}. Berikan \omterm{def:g11:gene-expression:protein}{proteinnya} dan golongkan.
        Substitusi tunggal mana yang
        menghasilkannya?
  \item Mutan 4: sebuah T disisipkan setelah \omterm{def:g10:universal-dna:nucleotide}{nukleotida} keenam:
        \texttt{ATG GCT TTG GAA ACC
        CGA ATA CGG TCA TTT CTG A}. Berikan \omterm{def:g11:gene-expression:protein}{proteinnya} dan golongkan.
  \item Urutkan keempat mutannya dari yang paling sedikit sampai yang
        paling banyak mengganggu \omterm{def:g11:gene-expression:protein}{proteinnya} dan berikan alasannya.
\end{enumerate}

\medskip
\noindent\textbf{Bagian III --- Mengapa kodenya seperti itu.}
\begin{enumerate}[resume]
  \item Berapa \omterm{def:g11:gene-expression:code}{kodon} yang menetapkan leusin? Serin? Triptofan? \omterm{def:g10:chemistry-of-life:families}{Asam
        amino} mana yang paling mungkin diubah oleh substitusi acak pada
        \omterm{def:g11:gene-expression:code}{kodonnya}, dan mana yang paling tidak mungkin?
  \item Hitung bagian seluruh substitusi pada kedudukan ketiga sebuah
        \omterm{def:g11:gene-expression:code}{kodon} leusin \texttt{CUx} yang bersifat diam.
  \item Jelaskan mengapa kode berhuruf dua tidak dapat bekerja, dan
        mengapa kode berhuruf empat tidak diperlukan.
  \item Poli-U memberi poli-fenilalanin dan poli-A memberi poli-lisin.
        Entri tabel mana yang ditegakkan kedua percobaan itu?
  \item Kodenya sama pada bakteri dan pada manusia. Nyatakan akibatnya
        bagi \omterm{prop:g10:universal-dna:universal}{transgenesis} dan akibatnya bagi kekerabatan.
\end{enumerate}

\medskip
\noindent\textbf{Bagian IV --- Mengukur waktu pabriknya.}
Sebuah \omterm{def:g10:cells-common-unit:organelle}{ribosom} menambahkan 10 \omterm{def:g10:chemistry-of-life:families}{asam amino} tiap sekon dan \omterm{def:g10:cells-common-unit:organelle}{ribosom} saling
menyusul berjarak 80 \omterm{def:g10:universal-dna:nucleotide}{nukleotida} di sepanjang sebuah mRNA; tiap mRNA
hidup 2 jam.
\begin{enumerate}[resume]
  \item Berapa lama satu \omterm{def:g10:cells-common-unit:organelle}{ribosom} menerjemahkan \omterm{def:g11:gene-expression:protein}{protein} \omterm{def:g10:universal-dna:gene}{gen} ini? Dan
        \omterm{def:g11:gene-expression:protein}{protein} berisi 600 \omterm{def:g10:chemistry-of-life:families}{asam amino}?
  \item Berapa banyak \omterm{def:g10:cells-common-unit:organelle}{ribosom} yang dapat membaca mRNA sepanjang 1800
        \omterm{def:g10:universal-dna:nucleotide}{nukleotida} sekaligus?
  \item Jika sebuah \omterm{def:g10:cells-common-unit:organelle}{ribosom} baru berangkat tiap 8 sekon, berapa salinan
        \omterm{def:g11:gene-expression:protein}{proteinnya} yang dihasilkan satu mRNA seumur hidupnya?
  \item Sebuah bakteri memerlukan \num{2000} salinan sebuah \omterm{def:g10:cell-metabolism:metabolism}{enzim} dalam
        10 menit. Berapa mRNA \omterm{def:g10:universal-dna:gene}{gennya}, yang ditranskripsikan sekaligus,
        yang diperlukan?
  \item Nyatakan hasilnya: \omterm{def:g11:gene-expression:protein}{protein} yang disandikan \omterm{def:g10:universal-dna:gene}{gennya}, \omterm{def:g11:mutations:mutation}{mutasi} di
        antara keempatnya yang memusnahkannya, dan waktu yang
        diperlukan satu \omterm{def:g10:cells-common-unit:organelle}{ribosom} untuk membangunnya.
\end{enumerate}
\end{problem}
